% chapter: 付録（数学の広がり）

\chapter{付録}
\begin{multicols*}{2}

\section{この付録で見るもの}

本編では,漸化式を与えられたときに,その一般項を求めるための考え方を整理してきた.
第1章から始まり,第2章では基本となる知識を,第3章では,漸化式を分類しそれぞれに対してのアプローチを扱った.
第4章ではそれらをまとめ上げ,本書の主題である"体系的解説"を行った.

付録では,問題を解くためのパターンを追加するのではなく,これまでの解法がより高度な数学で,あるいは別の視点から見て,どのような理論を内包しているかを眺める.
一つの漸化式を少し遠くから見ると,次のような問いが現れる.

\begin{itemize}
  \item 隣り合う項の差を,微分や積分に似た操作として扱えないか.
  \item $a_n$を無限に遠くまで見通すと最終的にどうなるのか.
  \item $a_n$の$n$を実数に拡張するとどうなるのか.
  \item 度々登場した特性方程式と不動点をより統一的に扱えないか.
\end{itemize}

これらの疑問をより広い数学の世界から見渡して,そこから生まれる新たな数学の道具が生まれる流れを追うものである.
その道具はしばしば大学数学を前提としたものになる.
しかし,付録を読むために,すべての大学数学を先に知っている必要はない.
あくまで順番に知的好奇心を満たすために,疑問を解決していくだけである.

各節では,まず問題意識を示し,次に必要な用語を定義し,最後に漸化式との
つながりを確認する.証明を完全に一般化することよりも,なぜその概念が
必要になり,何を見えるようにしたのかを大切にする.
% section: 差分から始まる離散的な微積分

\section{差分から始まる離散的な微積分}

\subsection{数列の変化を測る}

数列
\[
  a_0,a_1,a_2,\ldots
\]
の値だけを眺めていると,変化の規則が見えにくいことがある.
そこで,隣り合う二項の差を考える.
\[
  \Delta a_n=a_{n+1}-a_n
\]
これを数列の一階差分という.

例えば,
\[
  a_n=n^2
\]
ならば,
\[
  \Delta a_n=(n+1)^2-n^2=2n+1
\]
である.
さらに差分を取ると,
\[
  \Delta^2a_n
  =\Delta(2n+1)
  =(2(n+1)+1)-(2n+1)=2
\]
となる.

二次式が二回の差分で定数になった.
これは,二次関数を二回微分すると定数になることに対応している.
差分は,離散的な世界で変化の速さを測る操作なのである.

\subsection{差分を取ると次数が下がる理由}

多項式
\[
  p(n)=c_mn^m+c_{m-1}n^{m-1}+\cdots+c_0
\]
を考える.
二項定理から,
\[
  (n+1)^m-n^m
  =mn^{m-1}+\text{それより低い次数の項}
\]
である.
したがって,$\Delta p(n)$ の次数は,$p(n)$ の次数より一つ低い.

特に,$m$ 次多項式に $m$ 回差分を取ると定数になり,$m+1$ 回差分を取ると $0$ になる.
この事実から,数列が多項式で表されるかどうかを,差分表から推測できる.

\[
  \begin{array}{c|cccccc}
    n           & 0 & 1 & 2 & 3 & 4  & 5  \\ \hline
    a_n         & 0 & 1 & 4 & 9 & 16 & 25 \\
    \Delta a_n  & 1 & 3 & 5 & 7 & 9  &    \\
    \Delta^2a_n & 2 & 2 & 2 & 2 &    &
  \end{array}
\]

この表は,値の列よりも,変化の列の方が単純になることを示している.
漸化式を変形するという本編の発想も,実は「単純な変化を探す」という作業だった.

\subsection{差分演算子とシフト演算子}

差分の仕組みを,記号で整理してみよう.
数列を一つの対象と考え,シフト演算子 $E$ を
\[
  (Ea)_n=a_{n+1}
\]
で定める.
すると,
\[
  \Delta a_n=a_{n+1}-a_n=(E-1)a_n
\]
であるから,
\[
  \Delta=E-1
\]
と書ける.

この記号の利点は,差分を何回も取る操作を展開できることである.
\[
  \Delta^2=(E-1)^2=E^2-2E+1
\]
だから,
\[
  \Delta^2a_n=a_{n+2}-2a_{n+1}+a_n
\]
となる.
さらに,
\[
  \Delta^k=(E-1)^k
  =\sum_{j=0}^{k}(-1)^{k-j}\binom{k}{j}E^j
\]
より,
\[
  \Delta^ka_n
  =\sum_{j=0}^{k}(-1)^{k-j}\binom{k}{j}a_{n+j}
\]
を得る.

ここで現れた二項係数は偶然ではない.
差分は「一つ先へ進める操作」と「何もしない操作」の差であり,何度も差を取ると二項定理が現れるのである.

\subsection{差分の逆操作としての和}

微分の逆操作が積分であるように,差分の逆操作は和である.
例えば,
\[
  \Delta a_n=b_n
\]
ならば,
\[
  a_{n+1}-a_n=b_n
\]
なので,$0$ から $n-1$ まで足し合わせると,
\[
  a_n-a_0=\sum_{k=0}^{n-1}b_k
\]
となる.
従って,
\[
  a_n=a_0+\sum_{k=0}^{n-1}b_k
\]
である.

これは階差型漸化式の解法そのものである.
本編では解法として使ったが,ここでは「差分を元に戻す操作」として理解できる.

例えば,
\[
  a_{n+1}-a_n=2n+1,\qquad a_0=0
\]
ならば,
\[
  a_n=\sum_{k=0}^{n-1}(2k+1)
\]
であり,これは
\[
  a_n=n^2
\]
を与える.
平方数は,奇数を順番に足してできるという見方も得られた.

\subsection{和の公式を差分で設計する}

逆に,和
\[
  1^2+2^2+\cdots+n^2
\]
を求めたいとする.
公式を覚えていなくても,差分を利用して公式の形を設計できる.

和を $S_n$ とおくと,
\[
  S_n-S_{n-1}=n^2
\]
である.
右辺が二次式だから,$S_n$ は三次式になると予想する.
\[
  S_n=An^3+Bn^2+Cn+D
\]
と置けば,
\[
  \begin{aligned}
    S_n-S_{n-1}
     & =
    A\{n^3-(n-1)^3\}
    +B\{n^2-(n-1)^2\}
    +C   \\
     & =
    3An^2+(-3A+2B)n+(A-B+C)
  \end{aligned}
\]
である.
これが $n^2$ に一致するためには,
\[
  3A=1,\qquad -3A+2B=0,\qquad A-B+C=0
\]
である.
さらに $S_0=0$ より $D=0$ だから,
\[
  A=\frac13,\qquad B=\frac12,\qquad C=\frac16
\]
となる.
従って,
\[
  1^2+2^2+\cdots+n^2
  =\frac{n(n+1)(2n+1)}6
\]
が得られる.

ここで重要なのは,和の公式を暗記したことではない.
「差分を取ると元の被加数になるような式を探す」という,離散的な積分の発想を使ったことである.

% section: 反復としての漸化式と力学系

\section{反復としての漸化式と力学系}

\subsection{数列を点の運動として見る}

一項前から次の項を決める漸化式
\[
  a_{n+1}=f(a_n)
\]
を考える.
初期値 $a_0$ を決めると,
\[
  a_0,\ f(a_0),\ f(f(a_0)),\ldots
\]
という列ができる.

この列を,関数 $f$ によって動かされる点の軌道と見る.
数学では,このように状態の更新規則を調べる分野を\textbf{力学系}という.
物理の運動を扱わなくても,値が次々に変化するなら力学系として考えられる.

この視点では,一般項よりも次の問いが中心になる.

\begin{itemize}
  \item 軌道はある値へ近づくか.
  \item ある範囲から飛び出すか.
  \item 周期的に繰り返すか.
  \item 初期値を少し変えると,軌道は大きく変わるか.
\end{itemize}

\subsection{不動点と周期点}

$f(L)=L$ を満たす点 $L$ を不動点という.
不動点から出発した軌道は,
\[
  L\longmapsto L\longmapsto L\longmapsto\cdots
\]
と動かない.

一方,$f^{\,p}(L)=L$ を満たし,$p$ 回より少ない反復では元に戻らない点を周期 $p$ の周期点という.
例えば,
\[
  f(x)=2-x
\]
では $1$ が不動点であり,$x\neq1$ ならば
\[
  x\longmapsto 2-x\longmapsto x
\]
となるので,周期 $2$ の軌道になる.

「不動点を探す」だけでは,数列の全体の動きは分からない.
不動点へ近づく軌道もあれば,そこを飛び越えて周期軌道になるものもある.

\subsection{一次関数による反復を完全に分類する}

最も基本的な反復として,
\[
  a_{n+1}=ra_n+b
\]
を考える.
$r\neq1$ のときの不動点は,
\[
  L=rL+b
\]
より,
\[
  L=\frac{b}{1-r}
\]
である.
そこで $u_n=a_n-L$ と置くと,
\[
  \begin{aligned}
    u_{n+1}
     & =a_{n+1}-L \\
     & =ra_n+b-L  \\
     & =r(a_n-L)  \\
     & =ru_n
  \end{aligned}
\]
となる.
つまり,不動点からのずれは等比数列になる.
\[
  a_n-L=r^n(a_0-L)
\]
である.

この一行から,動きが完全に分かる.

\begin{itemize}
  \item $|r|<1$ なら $a_n\to L$.
  \item $r>1$ なら,一般には不動点から遠ざかる.
  \item $r<-1$ なら,符号を交互に変えながら大きくなる.
  \item $r=-1$ なら,通常は周期 $2$ になる.
  \item $r=1$ なら,$a_{n+1}=a_n+b$ であり,等差数列になる.
\end{itemize}

ここでは,漸化式を解く方法と,反復の動きを分類する方法が一致した.
一般項を求めることが,動力学的な理解へ直結している例である.

\subsection{微分は不動点の安定性を測る}

一般の関数 $f$ について,不動点 $L$ の近くで $a_n=L+u_n$ と置く.
テイラー展開の一次近似から,
\[
  f(L+u_n)
  \approx f(L)+f'(L)u_n
  =L+f'(L)u_n
\]
となる.
したがって,
\[
  u_{n+1}\approx f'(L)u_n
\]
である.

不動点の近くでは,結局,一次関数の反復に近い.
そこで,
\[
  |f'(L)|<1
\]
ならばずれが縮み,$L$ は安定であると考えられる.
逆に,
\[
  |f'(L)|>1
\]
ならばずれが拡大し,$L$ は不安定になりやすい.

この判定は,一般項を求めなくても数列の極限を予測するための強力な道具になる.
ただし,これは $L$ の近くでの局所的な判定である.
初期値が遠い場合には,別の不動点へ行ったり,発散したりする可能性がある.

\subsection{ロジスティック写像：単純な規則から複雑な動きへ}

次の漸化式を考える.
\[
  a_{n+1}=r a_n(1-a_n)
\]
これは,人口の割合や資源の制限を簡単に表すモデルとして現れる.
$0<a_n<1$ なら,$a_n$ は割合として解釈できる.

右辺の $ra_n$ は,現在の人口に比例する増加を表す.
一方,$(1-a_n)$ は,人口が大きくなるほど増加の余地が小さくなることを表している.
一つの式の中に,増加と抑制が同時に入っている.

不動点を求めると,
\[
  L=rL(1-L)
\]
より,
\[
  L=0,\qquad L=1-\frac1r
\]
である.
後者における微分係数は,
\[
  f'(L)=r(1-2L)=2-r
\]
となる.
したがって,
\[
  |2-r|<1
\]
すなわち
\[
  1<r<3
\]
の範囲では,正の不動点が安定になると予想される.

ところが $r$ をさらに大きくすると,不動点へ収束せず,周期 $2$,周期 $4$,さらに複雑な軌道が現れる.
この現象を完全に説明するには,力学系の理論が必要になる.
しかし,ここで知っておきたいのは,非常に単純な一行の漸化式でも,一般項を求めるという問題を超えて,安定性や周期性という深い問題を生み出すということである.

% section: 離散と連続を結ぶ

\section{離散と連続を結ぶ}

\subsection{時間の幅を明示する}

微分方程式
\[
  \frac{dx}{dt}=F(x)
\]
を考える.
時間を幅 $h$ で区切り,
\[
  t_n=nh,\qquad a_n=x(t_n)
\]
と置く.
微分の定義から,
\[
  \frac{x(t_{n+1})-x(t_n)}h
  \approx \frac{dx}{dt}(t_n)
  =F(x(t_n))
\]
である.
従って,
\[
  a_{n+1}=a_n+hF(a_n)
\]
という漸化式が得られる.

$h$ は単なる記号ではない.
それは,一回の更新がどれくらいの時間を表しているかを決める.
もし $h=1$ と暗黙に置くなら,そのことを意識しておく必要がある.

\subsection{オイラー法の誤差を一歩だけ調べる}

テイラー展開により,
\[
  x(t+h)
  =x(t)+hx'(t)+\frac{h^2}{2}x''(t)+\cdots
\]
である.
オイラー法は,
\[
  x(t+h)\approx x(t)+hx'(t)
\]
と,二次以上の項を捨てている.
従って,一回の更新で生じる誤差は,通常 $h^2$ 程度である.

しかし,時間 $T$ まで進むには,およそ $T/h$ 回の更新が必要である.
一回あたり $h^2$ の誤差が累積すると,全体の誤差はおおよそ
\[
  \frac{T}{h}\cdot h^2=Th
\]
程度になる.
これが,オイラー法の全体誤差が $h$ の一次程度になる理由である.

厳密な誤差評価には条件が必要だが,細かい刻みを使うほど近似がよくなる理由はこのように説明できる.

\subsection{指数関数は連続な等比数列である}

最も基本的な微分方程式
\[
  \frac{dx}{dt}=\lambda x
\]
を考える.
オイラー法では,
\[
  a_{n+1}=a_n+h\lambda a_n
  =(1+h\lambda)a_n
\]
となる.
これは等比数列である.
\[
  a_n=(1+h\lambda)^n a_0
\]

一方,微分方程式の解は,
\[
  x(t)=e^{\lambda t}x(0)
\]
である.
$t=nh$ とすれば,
\[
  e^{\lambda nh}=\left(e^{\lambda h}\right)^n
\]
だから,連続的な指数関数も,細かい時間幅で見れば等比的な更新の極限として現れる.

実際,
\[
  e^{\lambda h}
  =1+\lambda h+\frac{\lambda^2h^2}{2}+\cdots
\]
なので,$h$ が小さいとき,
\[
  1+\lambda h
\]
は $e^{\lambda h}$ の一次近似である.

「等比数列」と「指数関数」は別々の公式ではない.
一段ずつ進む世界と,連続的に進む世界で,同じ増加の仕組みを記述している.

\subsection{離散的な安定性と刻み幅}

$\lambda<0$ のとき,連続解 $e^{\lambda t}$ は $0$ へ近づく.
オイラー法の漸化式
\[
  a_{n+1}=(1+h\lambda)a_n
\]
が $0$ へ近づくためには,
\[
  |1+h\lambda|<1
\]
が必要である.
従って,
\[
  0<h<-\frac2\lambda
\]
でなければならない.

元の微分方程式が安定でも,刻み幅を大きく取りすぎると,近似の漸化式は発散する.
これは,「連続化したから安心」ではないことを示している.
離散化は新しい漸化式を作る操作であり,その漸化式自身の性質を調べなければならない.

% section: 解を直接求めず,近づく数列を設計する

\section{解を直接求めず,近づく数列を設計する}

\subsection{テイラー展開は何をしているのか}

関数 $f(x)$ を $x=a$ の近くで調べる.
接線による近似は,
\[
  f(x)\approx f(a)+f'(a)(x-a)
\]
である.
二次の項まで含めれば,
\[
  f(x)
  \approx f(a)+f'(a)(x-a)
  +\frac{f''(a)}{2}(x-a)^2
\]
となる.

テイラー展開は,関数を多項式で近似する道具である.
ただし,「近似だから適当でよい」という意味ではない.
どの点の近くで,何次まで使うかを決めることで,関数の局所的な情報を数値計算に変換している.

\subsection{ニュートン法の導出}

方程式
\[
  f(x)=0
\]
の解を探す.
現在の近似値を $x_n$ とし,$x_n$ の近くでは関数を接線で置き換える.
\[
  0\approx f(x_n)+f'(x_n)(x-x_n)
\]
を $x$ について解くと,
\[
  x=x_n-\frac{f(x_n)}{f'(x_n)}
\]
である.
この値を次の近似値に採用する.
\[
  x_{n+1}=x_n-\frac{f(x_n)}{f'(x_n)}
\]

ニュートン法は,方程式を解く公式ではない.
方程式を,収束することを期待した漸化式へ変換する設計法である.

\subsection{平方根の例と二次収束}

\[
  f(x)=x^2-2
\]
にニュートン法を使うと,
\[
  x_{n+1}=\frac{x_n^2+2}{2x_n}
\]
となる.
$x_0=1$ とすれば,
\[
  x_1=\frac32,\quad
  x_2=\frac{17}{12},\quad
  x_3=\frac{577}{408}
\]
である.

$\alpha=\sqrt2$ とおくと,
\[
  \begin{aligned}
    x_{n+1}-\alpha
     & =\frac{x_n^2+2}{2x_n}-\alpha             \\
     & =\frac{x_n^2-2\alpha x_n+\alpha^2}{2x_n} \\
     & =\frac{(x_n-\alpha)^2}{2x_n}.
  \end{aligned}
\]

誤差を $e_n=x_n-\alpha$ と書けば,
\[
  e_{n+1}=\frac{e_n^2}{2x_n}
\]
である.
誤差が二乗されるため,十分近くまで来ると正しい桁数が急速に増える.
この性質を二次収束という.

\subsection{一般の単純な根の近くで起こること}

$f(\alpha)=0$ かつ $f'(\alpha)\neq0$ とする.
$x_n=\alpha+e_n$ とおき,二次までテイラー展開すると,
\[
  f(x_n)=f'(\alpha)e_n+\frac{f''(\alpha)}2e_n^2+\cdots
\]
である.
同様に,
\[
  f'(x_n)=f'(\alpha)+f''(\alpha)e_n+\cdots
\]
となる.
これをニュートン法へ代入すると,一次の誤差が打ち消され,

\[
  e_{n+1}=C e_n^2+\text{三次以上の項}
\]
という形になる.
これがニュートン法が,根の近くで二次収束する理由である.

ただし,根が重なっていて $f'(\alpha)=0$ になる場合や,初期値が遠い場合には,この説明はそのまま使えない.
ニュートン法を適用する前に,何を近似しているのか,どの範囲で近似が有効なのかを確認する必要がある.

\subsection{微分を使わないニュートン法}

微分が計算しにくい場合,二つの近似値を使って接線の傾きを近似することができる.
\[
  f'(x_n)\approx
  \frac{f(x_n)-f(x_{n-1})}{x_n-x_{n-1}}
\]
と置くと,
\[
  x_{n+1}
  =
  x_n
  -
  f(x_n)
  \frac{x_n-x_{n-1}}{f(x_n)-f(x_{n-1})}
\]
を得る.
これは割線法と呼ばれる.

ニュートン法と割線法の違いは,単なる公式の違いではない.
ニュートン法は関数の傾きを正確に使い,割線法は過去の値から傾きを推測する.
どちらも「次の値をどう設計すれば,解へ近づくか」という同じ問題への答えである.

\input{第5章 付録/06_線形漸化式と線形代数/線形漸化式と線形代数.tex}
% section: チェビシェフ多項式：三角関数を多項式に翻訳する

\section{チェビシェフ多項式：三角関数を多項式に翻訳する}

\subsection{同じ漸化式が二つの姿を持つ}

\[
  a_{n+1}=2x a_n-a_{n-1}
\]
を考える.
$x$ を固定すれば二階線形漸化式であり,特性方程式は
\[
  r^2-2xr+1=0
\]
である.

$-1\leq x\leq1$ のとき,$x=\cos\theta$ と書ける.
すると,
\[
  \cos((n+1)\theta)
  =
  2\cos\theta\cos(n\theta)-\cos((n-1)\theta)
\]
だから,$\cos(n\theta)$ は同じ漸化式を満たす.

ここで,「三角関数の値を計算する問題」と「漸化式を解く問題」が同じ式を持っている.
この一致を,変数 $x$ の多項式として保存したくなる.

\subsection{チェビシェフ多項式の定義と証明}

第一種チェビシェフ多項式 $T_n(x)$ を,
\[
  T_0(x)=1,\qquad T_1(x)=x
\]
および
\[
  T_{n+1}(x)=2xT_n(x)-T_{n-1}(x)
\]
で定める.

初めのいくつかは,
\[
  \begin{aligned}
    T_0(x) & =1,          \\
    T_1(x) & =x,          \\
    T_2(x) & =2x^2-1,     \\
    T_3(x) & =4x^3-3x,    \\
    T_4(x) & =8x^4-8x^2+1
  \end{aligned}
\]
である.

漸化式の右辺が多項式だけでできているので,すべての $T_n(x)$ は多項式になる.
一方,
\[
  T_n(\cos\theta)=\cos(n\theta)
\]
が成り立つ.

実際,$n=0,1$ では定義から明らかである.
$n$ と $n-1$ について成り立つと仮定すれば,
\[
  \begin{aligned}
    T_{n+1}(\cos\theta)
     & =2\cos\theta T_n(\cos\theta)-T_{n-1}(\cos\theta) \\
     & =2\cos\theta\cos(n\theta)-\cos((n-1)\theta)      \\
     & =\cos((n+1)\theta)
  \end{aligned}
\]
となる.

三角関数を使って書かれた量を,$x$ の多項式に翻訳できた.
例えば,
\[
  \cos(3\theta)=4\cos^3\theta-3\cos\theta
\]
は,
\[
  T_3(x)=4x^3-3x
\]
という多項式の公式になる.

\subsection{三種類の領域：振動,境界,増大}

特性方程式
\[
  r^2-2xr+1=0
\]
の根は,
\[
  r=x\pm\sqrt{x^2-1}
\]
である.

\begin{itemize}
  \item $|x|<1$ なら根は複素数で,単位円上にある.$T_n(x)$ は振動する.
  \item $x=1$ または $x=-1$ では根が重なり,境界的な振る舞いになる.
  \item $|x|>1$ なら実根の絶対値が異なり,一方の根が支配的になる.
\end{itemize}

$x=\cos\theta$ の代わりに $x=\cosh t$ と置くと,
\[
  T_n(\cosh t)=\cosh(nt)
\]
となる.
三角関数の振動と双曲線関数の増大が,同じ多項式の異なる領域で現れている.

これは,特性根の絶対値が数列の増加を支配するという,線形漸化式の基本原理を別の言葉で見たものである.

\subsection{第二種チェビシェフ多項式}

第一種だけでなく,\[
  U_0(x)=1,\qquad U_1(x)=2x
\]
および
\[
  U_{n+1}(x)=2xU_n(x)-U_{n-1}(x)
\]
で定まる第二種チェビシェフ多項式も考えられる.
こちらは,
\[
  U_n(\cos\theta)
  =\frac{\sin((n+1)\theta)}{\sin\theta}
\]
を満たす.

第一種が余弦を記録するのに対して,第二種は正弦を記録する.
同じ漸化式でも初期値を変えると,異なる多項式族が現れる.
初期値は単なる付属条件ではなく,解の空間の中からどの解を選ぶかを決めている.

\subsection{直交性と近似への入口}

チェビシェフ多項式が特別なのは,三角関数を表すからだけではない.
余弦の直交性から,
\[
  \int_0^\pi\cos(m\theta)\cos(n\theta)\,d\theta=0
  \qquad(m\neq n)
\]
が得られる.

$x=\cos\theta$ と変数変換すれば,
\[
  \int_{-1}^{1}
  \frac{T_m(x)T_n(x)}{\sqrt{1-x^2}}\,dx=0
  \qquad(m\neq n)
\]
となる.
この性質を直交性という.

直交する多項式を使うと,複雑な関数を
\[
  c_0T_0(x)+c_1T_1(x)+c_2T_2(x)+\cdots
\]
のように展開して近似できる.
ここからフーリエ級数や数値計算,最良近似の理論へ進むことができる.

本書で必要なのは近似理論の完成ではない.
一つの漸化式が,三角関数を経由して,多項式の直交性や関数近似へつながるという道筋を見つけることである.

\clearpage

\input{第5章 付録/08_特殊関数の広がり/特殊関数の広がり.tex}
% section: 母関数：無限個の数を一つの対象にする

\section{母関数：無限個の数を一つの対象にする}

\subsection{なぜ数列を式に詰め込むのか}

数列は無限個の項を持つ.
項を一つずつ扱う代わりに,すべての項を係数として一つの式にまとめる.
\[
  A(x)=a_0+a_1x+a_2x^2+\cdots
\]
これを母関数という.

母関数では,$x^n$ の係数が $a_n$ である.
したがって,母関数を変形して係数を読むことができれば,数列の問題を式の問題へ移せる.

\subsection{形式的冪級数としての母関数}

ここで,母関数を通常の関数と考える必要はない.
例えば,
\[
  1+x+x^2+x^3+\cdots
\]
は,$|x|<1$ では $1/(1-x)$ に収束するが,形式的冪級数としては,係数を並べた無限列である.

形式的に,
\[
  (1-x)(1+x+x^2+\cdots)=1
\]
と係数を比較できる.
収束するかどうかを考えなくても,漸化式の係数操作が可能になる.

一方,実際に数値を代入して関数として使う場合には,収束範囲を確認する必要がある.
同じ記号でも,形式的な議論と解析的な議論を区別しなければならない.

\subsection{線形漸化式を母関数へ移す}

\[
  a_{n+2}=a_{n+1}+a_n
\]
とし,
\[
  A(x)=\sum_{n=0}^{\infty}a_nx^n
\]
と置く.
漸化式に $x^n$ を掛けて足すと,
\[
  \sum_{n=0}^{\infty}a_{n+2}x^n
  =
  \sum_{n=0}^{\infty}a_{n+1}x^n
  +
  \sum_{n=0}^{\infty}a_nx^n
\]
である.

左辺と右辺を $A(x)$ で表すと,
\[
  \frac{A(x)-a_0-a_1x}{x^2}
  =
  \frac{A(x)-a_0}{x}+A(x)
\]
となる.
これを整理して,
\[
  A(x)=\frac{a_0+(a_1-a_0)x}{1-x-x^2}
\]
を得る.

分母の多項式 $1-x-x^2$ は,漸化式の特性方程式と同じ情報を持つ.
母関数は,特性方程式を別の形で記録しているのである.

\subsection{畳み込みが積になる}

母関数の特に重要な性質は,積の係数である.
\[
  A(x)=\sum_{n=0}^{\infty}a_nx^n,\qquad
  B(x)=\sum_{n=0}^{\infty}b_nx^n
\]
とすると,
\[
  A(x)B(x)
  =
  \sum_{n=0}^{\infty}
  \left(\sum_{k=0}^{n}a_kb_{n-k}\right)x^n
\]
となる.

内側の和
\[
  \sum_{k=0}^{n}a_kb_{n-k}
\]
を畳み込みという.
二つの選択を組み合わせる数え上げでは,畳み込みが自然に現れる.
母関数は,畳み込みを積へ変えることで,組合せ論の漸化式を代数方程式へ移す.

\subsection{カタラン数：非線形漸化式を解く}

カタラン数 $C_n$ を,
\[
  C_0=1,\qquad
  C_{n+1}=\sum_{k=0}^{n}C_kC_{n-k}
\]
で定める.
最初の項は,
\[
  1,1,2,5,14,42,\ldots
\]
である.

母関数
\[
  C(x)=\sum_{n=0}^{\infty}C_nx^n
\]
を考える.
漸化式の右辺は畳み込みなので,母関数では積になる.
\[
  C(x)-1=xC(x)^2
\]
すなわち,
\[
  xC(x)^2-C(x)+1=0
\]
である.

二次方程式として解くと,
\[
  C(x)=\frac{1\pm\sqrt{1-4x}}{2x}
\]
となる.
$x=0$ の近くで $C(x)\to C_0=1$ でなければならないので,正しい枝は,
\[
  C(x)=\frac{1-\sqrt{1-4x}}{2x}
\]
である.

もともとの漸化式は,各項がそれ以前の項の積の和で表される非線形なものだった.
それが,母関数を導入すると二次方程式になった.
母関数は線形漸化式専用の道具ではない.
「組合せの分解が畳み込みになる」場面では,非線形な漸化式にも力を発揮する.

\subsection{母関数から見る漸化式の意味}

母関数による変換は,
\[
  \text{無限個の数}
  \longrightarrow
  \text{一つの形式的な対象}
\]
である.
この変換によって,シフトは $x$ の掛け算になり,畳み込みは積になり,漸化式は代数方程式になる.

何を簡単にしたいかによって,どの対象へ移すかが決まる.
差分では変化を簡単にし,行列では連立した状態を簡単にし,母関数では無限個の項の関係を簡単にする.

% section: 分数一次変換と離散リカッチ方程式

\section{分数一次変換と離散リカッチ方程式}

\subsection{一次分数型は関数の反復である}

\[
  a_{n+1}=\frac{pa_n+q}{ra_n+s}
\]
を考える.
これは一次分数変換,またはメビウス変換と呼ばれる.
分母があるため一見扱いにくいが,関数の反復として見ると構造が現れる.

不動点は,
\[
  \alpha=\frac{p\alpha+q}{r\alpha+s}
\]
より,
\[
  r\alpha^2+(s-p)\alpha-q=0
\]
を満たす.
従って,不動点は二次方程式の解として求められる.

\subsection{二つの不動点を基準にする}

異なる二つの不動点を $\alpha,\beta$ とする.
次の変数変換を考える.
\[
  b_n=\frac{a_n-\alpha}{a_n-\beta}
\]
これは,$a_n$ が $\alpha$ と $\beta$ のどちらに近いかを,比として記録する量である.

変換
\[
  f(x)=\frac{px+q}{rx+s}
\]
について,固定点の条件を使うと,
\[
  f(x)-\alpha
  =
  \frac{(p-r\alpha)(x-\alpha)}{rx+s}
\]
および
\[
  f(x)-\beta
  =
  \frac{(p-r\beta)(x-\beta)}{rx+s}
\]
となる.
従って,
\[
  \begin{aligned}
    \frac{f(x)-\alpha}{f(x)-\beta}
     & =
    \frac{p-r\alpha}{p-r\beta}
    \frac{x-\alpha}{x-\beta}
  \end{aligned}
\]

ここで
\[
  \lambda=\frac{p-r\alpha}{p-r\beta}
\]
と置けば,
\[
  b_{n+1}=\lambda b_n
\]
となる.

分数を含む複雑な漸化式が,変数変換によって等比数列になった.
本編で扱った一次分数型の解法は,このような「二つの基準点からの比を取る」という幾何学的な意味を持っている.

\subsection{連続リカッチ方程式との共通点}

連続的なリカッチ方程式では,既知の特別解を基準にして,
\[
  x=x_0+\frac1u
\]
と置いた.
離散的な一次分数型では,二つの不動点を基準にして,
\[
  b=\frac{x-\alpha}{x-\beta}
\]
と置いた.

変換の形は違うが,発想は同じである.
\[
  \text{非線形で扱いにくい変化}
  \longrightarrow
  \text{基準となる解や点との差}
  \longrightarrow
  \text{線形または等比的な変化}
\]

漸化式と微分方程式を結ぶのは,公式の対応ではない.
「どの量を記録すれば,変化が単純になるか」という問題意識の対応である.

\input{第5章 付録/11_漸近解析/漸近解析.tex}
\input{第5章 付録/12_漸化式が組合せ論と確率へ進む/漸化式が組合せ論と確率へ進む.tex}
\input{第5章 付録/13_漸化式と数論/漸化式と数論.tex}
\input{第5章 付録/14_解けない漸化式と解くことの意味/解けない漸化式と解くことの意味.tex}
\input{第5章 付録/15_付録の地図を閉じる/付録の地図を閉じる.tex}
\end{multicols*}
